\section{序論}

本章では，本研究が取り組む問題の背景を述べ，既存の取り組みに残された課題を整理したうえで，本研究の目的と本論文の構成について説明する．

\subsection{研究背景}

\subsubsection{大規模言語モデルの普及}

2022 年の ChatGPT の公開以降，大規模言語モデル（Large Language Model，以下 LLM）は急速に利用が広がった．文書の要約，翻訳，質問応答，プログラムの作成といった作業に用いられ，専門的な訓練を受けていない利用者が日常的に使う道具になっている．日本語環境でも，汎用モデルに加えて Swallow \cite{fujii2024swallow} のような日本語を主対象としたモデルが公開されており，日本語での評価基盤も整備されつつある \cite{llmjpeval2024}．

利用の広がりに伴い，LLM の出力をどこまで信じてよいかという問題が生じている．LLM は事実と異なる内容を生成することがあり，この現象は幻覚（hallucination）と呼ばれる \cite{ji2023hallucination,huang2023hallucination}．存在しない文献を引用する，起きていない出来事を詳細に説明する，誤った計算結果を提示するといった例が報告されている．

\subsubsection{語調が正しさを反映しない問題}

幻覚が問題になるのは，単に出力が誤っているからではない．より本質的なのは，\textbf{誤った出力が正しい出力と区別のつかない語調で提示される}ことである．LLM は正答時も誤答時も同じように流暢な文章を生成するため，利用者は出力の見た目から正しさを推し量れない．

もし LLM が「この回答には自信がない」と適切に伝えられれば，利用者はその回答を優先して検証できる．しかし現状では出力の語調からその手がかりが得られないため，どの回答に検証の手間を割くべきかを判断しにくい．出典の確認や重要度に応じた検証といった運用はありうるが，そこに検証の労力をどう配分するかの根拠が乏しい．

\subsubsection{確信度の表明と較正}

この問題への一つの対処として，LLM に自身の回答への確信度（confidence）を表明させる手法が研究されている \cite{kadavath2022know,lin2022teaching,tian2023justask}．具体的には「回答とともに，それが正しい確率を答えてください」といった指示をプロンプトに含め，確信度を言語として出力させる．この手法は verbalized confidence と呼ばれる．

表明された確信度が有用であるためには，それが実際の正答率と対応している必要がある．確信度 0.9 と答えた問題群の正答率が実際に 9 割程度であれば，利用者は「確信度 0.9 以上の回答は採用し，それ未満は検証する」という判断ができる．この対応の度合いを較正（calibration）と呼び，よく対応している状態をwell-calibrated という．較正の程度を測る代表的な指標が ECE（Expected Calibration Error）である \cite{guo2017calibration}．

\subsection{課題}

\subsubsection{英語以外の言語での検証が乏しい}

較正に関する研究の多くは英語のタスクを対象としている．LLM の性能が言語によって異なることは広く知られており，較正の性質も言語間で異なる可能性がある．多言語での較正を扱った研究として MlingConf \cite{xue2025mlingconf} があるが，扱われた言語は限られており，日本語のタスクに焦点を当てた体系的な検証は乏しい．

日本語での検証が必要な理由は二つある．第一に，日本語の利用者が LLM の確信度表示をどこまで信頼できるかという実用上の問いに，現状では答えられない．第二に，日本語特有の表現，とくに確信の度合いを表す言い回しが確信度の表明にどう影響するかは，英語の研究からは推し量れない．

\subsubsection{確信度の尋ね方による結果の変動}

確信度の引き出し方によって較正の結果が変わることは，複数の研究で指摘されている \cite{tian2023justask,xiong2024llmuncertainty}．回答と確信度を同時に尋ねるか，回答を得てから改めて尋ねるか，数値で答えさせるか言語表現で答えさせるかといった違いが，結果に影響する．

方式の比較そのものは Tian ら \cite{tian2023justask} をはじめとして英語では行われている．一方，日本語のタスクで，同一のモデル・同一のデータ・同一の評価指標のもとに複数の方式を並べた報告は限られている．方式ごとに異なる条件で得られた数値を並べても，差が方式に由来するのか条件の違いに由来するのかを切り分けられない．

\subsubsection{評価指標そのものの妥当性}

較正の評価には ECE が広く用いられているが，ECE の値は確信度の分布に左右される．確信度が狭い範囲に偏っている場合，多くの標本が少数の区間に入るため，ECE は全体の平均確信度と平均正答率の差に近い値をとりうる．このとき ECE は較正の水準を測ってはいるものの，その値だけから較正の良否と正答率の高低を読み分けることは難しい．本研究の予備実験でもこの状況が生じた（第 6 章）．評価指標の選択と併記の仕方は，較正研究の設計上の論点である．

\subsection{本研究の目的}

本研究は，日本語タスクにおいて，確信度の引き出し方の違いが LLM の較正に与える影響を明らかにすることを目的とする．そのうえで，較正を改善するプロンプト設計の指針を得ることを目指す．

具体的な問いは以下の通りである．

\begin{enumerate}
\item 日本語タスクにおいて，確信度の引き出し方式によって較正はどのように
変化するか．

\item 観察された差は，較正の水準の違いによるものか，それとも正答率の違いに
伴うものか．

\item 較正を改善するために，プロンプトをどのように設計すべきか．
\end{enumerate}

第 1 の問いに答えるため，3 つの引き出し方式を設計し，同一条件で比較する評価基盤を構築した．第 2 の問いについては，複数の指標を併記し，確信度の分布を提示することで解釈の材料を揃える．ただし本論文の範囲では，方式によって回答そのものが変わるため，較正の違いと正答率の違いを完全に分離することはできていない．分離には回答を固定した比較が必要であり，今後の課題とする．第 3 の問いは本論文の時点では未着手であり，予備実験から得られた知見をもとに今後取り組む．

本研究の成果は，日本語で LLM を利用する際に確信度表示をどこまで信頼できるかという判断材料を提供する．また，較正を測る際にどの指標を用いるべきかという方法論上の知見も提供する．

\subsection{本論文の構成}

本論文の構成は以下の通りである．第 2 章では較正の評価指標と確信度の表明に関する関連研究を概観し，本研究の位置づけを示す．第 3 章では確信度の引き出し方式，データセット，評価指標の設計方針を述べる．第 4 章では評価基盤の実装について述べる．第 5 章では予備実験の設定と結果を示す．第 6 章では結果を解釈し，本研究の限界を述べる．第 7 章では現時点のまとめと今後の方針を述べる．

\subsection{本章のまとめ}

本章では，LLM が誤った出力を正しい出力と同じ語調で提示するために，利用者が検証の労力を配分しにくいという問題を述べた．その対処として確信度を言語で表明させる手法があり，表明された確信度と実際の正答率の一致，すなわち較正が重要になる．既存研究は英語のタスクが中心であること，確信度の尋ね方によって結果が変わること，そして評価指標である ECE の値が確信度の分布に左右されることを課題として挙げた．これらを踏まえ，日本語タスクにおける引き出し方式の比較を本研究の目的として設定した．
